\chapter{Notation and a working glossary}
\label{app:notation}

The subject joins several mathematical languages. The same letter
can have a conventional meaning in more than one of them. This
appendix records the principal conventions and identifies the
distinctions that matter when moving between chapters.

\section{Fields and states}

\begin{longtable}{p{0.23\textwidth}p{0.70\textwidth}}
\toprule Symbol&Meaning\\\midrule
\endhead
$\one$&Vacuum state, tensor unit, or identity object according to
context. The vacuum simple-module label is also written $0$.\\
$Y(a,z)$&Field associated to state $a$;
$Y(a,z)=\sum_n a_{(n)}z^{-n-1}$.\\
$a_{(n)}b$&The coefficient of $(z-w)^{-n-1}$ in the OPE of $a$
and $b$; it has weight $\wt(a)+\wt(b)-n-1$.\\
$\mathsf T$, $D$&Translation operator on a VOA, equal to $L_{-1}$.
The construction chapters often use $D$.
It is distinct from a stress field $T(z)$ and a modular matrix $T$.\\
$\omega$, $T(z)$&Conformal vector and corresponding stress field;
$Y(\omega,z)=T(z)=\sum_nL_nz^{-n-2}$.\\
$H_L,H_R$&In the construction chapters, Cartan currents of the internal affine factors at
levels one and thirty-nine. Both are holomorphic; the subscripts
do not denote the two physical chiralities. In the chain chapter,
$H_L$ instead denotes the Hamiltonian of a chain of length $L$.\\
$\omega_L,\omega_R$&Sugawara conformal vectors of those internal
factors, with charges $1,117/41$ and fields $T_L,T_R$.\\
$a_{(-1)}b$&Normal-product state, corresponding to
$:Y(a,z)Y(b,z):$.\\
$h,\bar h$&Holomorphic and antiholomorphic conformal weights.\\
$\Delta$&Full scaling dimension $h+\bar h$.\\
$s_{\rm spin}$&Conformal spin $h-\bar h$. Some local physics
formulas abbreviate it to $s$; this is distinct from the
categorical abbreviation $1+d^2$.\\
$c$&One chiral central charge. For the nonchiral theories here,
$c_L=c_R=c$; their sum is $2c$.\\
$v$&The nonuniversal velocity in a lattice-to-cylinder energy
conversion; elsewhere a generic state vector when no velocity
is involved.\\
$\HH_{\mathrm{path},L}$&The length-$L$ periodic chain Hilbert
space, with orthonormal basis the admissible closed fusion paths.\\
$U_{\rm tr}$&The chain's unitary right cyclic shift, with
$U_{\rm tr}^L=1$ and $U_{\rm tr}h_jU_{\rm tr}^{-1}=h_{j+1}$.
It is distinct from the chiral algebra $U$ and VOA translation
$\mathsf T$.\\
$\widehat a$&A field divided by its positive norm. It does not
fix the remaining phase of a complex field.\\
$a^\dagger$&Radial adjoint or Hilbert-space adjoint, according
to context; not an instruction to add a right-moving field.\\
$\theta$&An antilinear PCT operation in VOA sections;
$\theta_i$ can also denote a categorical twist, distinguished
by its indexed scalar use.\\
\bottomrule
\end{longtable}

\section{Categories and modular data}

\begin{longtable}{p{0.23\textwidth}p{0.70\textwidth}}
\toprule Symbol&Meaning\\\midrule
\endhead
$\mathcal H_3$&The specified unitary six-object Haagerup
fusion category used for the physical chain.\\
$\mathcal H_2$&A different unitary category with the same
six-object fusion ring, Morita equivalent to $\mathcal H_3$.\\
$\alpha$&The invertible order-three Haagerup object. A lattice
root is also conventionally denoted by $\alpha$; its lattice
context or subscript distinguishes it.\\
$\rho$&A noninvertible Haagerup simple with
$\rho^2=\one+\rho+\alpha\rho+\alpha^2\rho$.\\
$d$&The positive quantum dimension
$(3+\sqrt{13})/2$, satisfying $d^2=1+3d$.\\
$s$&In the category and extension formulas,
$s=1+d^2=2+3d$.\\
$\operatorname{Dim}(\mathcal C)$&Global dimension,
$\sum_i d_i^2$, for a positive spherical fusion category.\\
$\mathcal D$&Total quantum dimension
$\sqrt{\operatorname{Dim}(\mathcal C)}$ of the modular category
under discussion. For $\mathcal C=Z(\mathcal H_3)$,
$\mathcal D=3s=\operatorname{Dim}(\mathcal H_3)$.\\
$Z(\mathcal C)$&Drinfeld center, whose objects have compatible
half-braidings with every object of $\mathcal C$.\\
$F^{abc}_d$&Associator matrix between two bases of a fusion
space. Left channels index rows and right channels index columns;
see \eqref{eq:F-basis-convention}.
Elsewhere $F$ can denote a lowering current or a forgetful
functor; its arguments distinguish these uses.\\
$X^*$&Categorical dual object. Dualizing a morphism is linear;
it differs from taking a Hermitian adjoint, which conjugates scalars.\\
$S$&Modular matrix, normalized by $S_{00}=1/\mathcal D$.
For a pointed lattice example we use the negative Fourier sign.\\
$\Theta$&Categorical twist matrix, with entries
$e^{2\pi ih_i}$.\\
$T_{\rm ch}$&Character transformation matrix,
$e^{-2\pi ic/24}\Theta$. At $c=8$ its common prefactor is
$e^{-2\pi i/3}$.\\
$a$&One of the twelve central simple labels. It is not the
invertible Haagerup object $\alpha$.\\
$\mathcal C_A$&All right modules of an algebra object $A$.
For the connected commutative separable algebras in the braided
setting used here, the braiding supplies a left action and
defines the balanced tensor product $\otimes_A$.\\
$\mathcal C_A^0$&Local modules: if $r:X\otimes A\to X$ is
the action, then $r\,c_{A,X}c_{X,A}=r$. This is an equation
for the action, not necessarily identity of the double
braiding on the entire underlying tensor product.\\
${}_A\mathcal C_A$&Algebra bimodules, describing the relevant
full-theory defects in the rational construction.\\
$\mathcal C^{\rm op}$&The convention in this book reverses
the order of the tensor product, not the direction of
morphisms.\\
$\mathcal C^{\rm rev}$&For a braided category, reverse the
braiding and twist. This differs from reversing tensor order.\\
$\mathcal C\boxtimes\mathcal D$, $M\boxtimes_V N$&The first is
the external product of two categories, with pairs of simple labels.
The second fuses two modules of one VOA $V$. Neither is the plain
tensor product of the underlying module vector spaces.\\
\bottomrule
\end{longtable}

\section{The main algebras and their charges}

\begin{longtable}{p{0.34\textwidth}p{0.59\textwidth}}
\toprule Object&Role\\\midrule
\endhead
$V_L$&Positive even lattice VOA, with central charge
$\rank L$.\\
$B=V_{E_6\oplus A_2}$&The immediate ambient algebra of the
six-current construction; $c=8$ and $\dim B_1=86$.\\
$H=V_{E_8}$&The original holomorphic lattice extension of
$B$; $c=8$ and $\dim H_1=248$.\\
$U$&The actual vertex algebra generated by the two affine
current triples, the ambient conformal vector, and the field
$u$. Under the finite closure inputs, $\dim U_1=6$.\\
$H'$&The holomorphic extension selected through the canonical
$\mathcal H_3$ algebra. It is isomorphic to $V_{E_8}$ but its
embedding of $U$ is not automatically the original one.\\
$A_{1,k}$&The simple integrable affine
$L_k(\mathfrak{sl}_2)$ VOA, of charge $3k/(k+2)$.\\
$K(\mathfrak{sl}_2,k)$&The parafermion commutant of the
Cartan Heisenberg algebra, of charge $2(k-1)/(k+2)$.\\
Potts&The $k=3$ parafermion algebra, of charge $4/5$.\\
Level-thirteen parafermion&The $k=13$ factor, of charge $8/5$.\\
$\mathcal T$&The rank-two lattice--Potts core, of charge
$14/5$. Its detailed lattice extension is defined in
Chapter~\ref{ch:factors}.\\
$A_L$&The larger order-thirteen core extension, of charge
$22/5$ and fourteen simple modules. The subscript is part of
the core notation, not a declaration that it is an affine
level.\\
$C_6$&The rational central-charge-six commutant constructed
from the saturated plane lattice under the accepted $c=8$
construction.\\
$c_*$&The irrational chiral charge
$2.025004294\ldots$ arising from a stress axis. It is not an
identified central charge of the physical chain.\\
\bottomrule
\end{longtable}

\section{A glossary of the main logical distinctions}

\begin{description}
\item[Abelian current.]A current whose OPE with the other currents
in the selected abelian family has only a central double pole
and no Lie-bracket simple pole. The double pole need not vanish.

\item[Anomaly.]An obstruction to a specified symmetry operation,
such as gauging a group action with nontrivial associator class.
A chiral gravitational anomaly and a finite-group associator
anomaly are different data; the context must specify which is meant.

\item[Associator.]A coherent change between the two ways of
parenthesizing a tensor product. Its pentagon relation is more
information than associativity of the fusion multiplicities.

\item[Boundary category.]The category of boundary conditions or
all extension modules in the specified construction. It can be
nonbraided even though the bulk chiral category is braided.

\item[Bounded localized vertex operators.]A property combining
bounded annulus-regularized insertion maps, with suitable energy
damping on the inserted state, and locality of the bounded
algebras they generate. It is stronger than boundedness of an
insertion for each fixed state separately.

\item[Branching.]Decomposition of a representation after restricting
to a subalgebra. Multiplicity spaces can carry a commuting algebra.

\item[Canonical.]A construction determined naturally after its
input data and conventions have been specified. The word does not
remove the need to state which category, algebra, or normalization
is being used.

\item[Cauchy complete category.]A category in which every
idempotent splits: if $p^2=p$ on $X$, there is an object
representing its image, with inclusion and projection composing
to $p$. This is the categorical version of retaining the range
of a projection as an object.

\item[CFT type.]A nonnegatively integer-graded VOA with
finite-dimensional homogeneous spaces and
$V_0=\C\one$. This does not assert rationality.

\item[Character.]A trace counting states by conformal energy.
It forgets most operator products and can coincide for distinct
simple modules.

\item[Commutant.]Fields whose singular OPE with a subalgebra
vanishes. Constructing a commutant is a chiral operation.

\item[Completely positive map.]A linear map preserving positive
semidefinite matrices even after tensoring with the identity on
any finite auxiliary matrix space. The MPS transfer formula
$\mathcal E(X)=\sum_sA^sX(A^s)^\dagger$ makes this property
explicit; trace preservation is a separate condition.

\item[Complete unitarity.]Positivity of module forms and canonical
fusion forms, beyond positivity of the vacuum space.

\item[Completely rational net.]A chiral conformal net satisfying
the split and strong-additivity properties and having finite
two-interval index. These hypotheses give a finite unitary
modular sector theory.

\item[Conformal block.]A chiral solution of the symmetry Ward
identities for specified external representations. A factorization
basis labels blocks by intermediate representations and junction
multiplicities, summing their descendants. A general block may be
a linear combination of that basis. Full correlators additionally
specify how holomorphic and antiholomorphic blocks are paired.

\item[Conformal net.]An assignment of a von Neumann algebra of
bounded observables to each interval, satisfying inclusion,
locality, conformal covariance, and the vacuum and positive-energy
conditions used in the text.

\item[Construction input.]A model-specific finite identity,
membership, rank, or structural datum accepted from the exact
construction record. It is a mathematical premise with a stated
verification scope.

\item[Core of an operator.]A subspace of an operator's domain
dense in its graph norm. Approximating a vector there also
approximates its image. Ordinary Hilbert-space density alone
does not supply this property.

\item[Effective central charge.]The parameter governing a leading
asymptotic trace or entropy under its stated spectral assumptions;
it need not equal the Virasoro central charge without a
normalizable vacuum of lowest weight zero.

\item[Energy bound.]An estimate controlling field modes by a
polynomial in their mode number and in the conformal Hamiltonian.
Such estimates make smearing and operator closure possible on
controlled domains.

\item[\'Etale algebra.]A commutative separable algebra object
in the braided category under discussion. Connectedness is
an additional condition, $\dim\operatorname{Hom}(\one,A)=1$.

\item[Exact finite result.]A statement about the specified finite
algebraic calculation or finite chain, proved with appropriate
arithmetic or certification. It does not alone control an
infinite-size limit.

\item[Fiber functor.]A faithful exact strong monoidal functor
from a fusion category to finite-dimensional vector spaces.
It is much more structure than a representation of the fusion
ring by ordinary matrices.

\item[Frobenius--Schur indicator.]The trace of pivotal cyclic
rotation on $\Hom(\one,X^{\otimes n})$. It detects information
about the associator and duality beyond the fusion rules. In a
spherical fusion category it also equals the categorical twist
trace of center induction divided by the global dimension.

\item[Fusion category.]A rigid semisimple complex linear tensor
category with finitely many simple isomorphism classes,
finite-dimensional morphism spaces, and simple tensor unit.
It need not have a braiding or a positive adjoint.

\item[Fusion space.]A junction multiplicity space such as
$\Hom(c,a\otimes b)$. Its dimension is a fusion coefficient,
not the dimension of a physical lattice Hilbert space.

\item[Gap.]An energy gap is an excitation-energy separation above
the vacuum sector. A scalar scaling gap is the lowest nonidentity
scalar dimension $\Delta$. A positive scaling gap in a CFT
corresponds to energies of order $1/L$.

\item[Haag duality.]On the conformal vacuum Hilbert space,
the algebra assigned to an interval equals the commutant of
the algebra assigned to its complementary interval. The
representation and cyclic-vacuum hypotheses matter.

\item[Holomorphic.]A rational chiral algebra with only its vacuum
module. The corresponding full theory still has infinitely many
local fields.

\item[Invariant field.]A field on which each simple symmetry line
acts by its vacuum eigenvalue, equal to that line's quantum
dimension.

\item[Ishibashi state.]A formal solution of chiral boundary
gluing equations within a compatible representation pair.
Its propagated overlap is a character; the unregularized
state is not generally normalizable.

\item[Isotropic subgroup.]A subgroup of a finite discriminant
form on which the quadratic function vanishes modulo one.
It determines an even lattice extension in the lattice setting.

\item[Jones--Wenzl projection.]The Temperley--Lieb idempotent
on a fixed number of parallel strands killed by the adjacent
turnback operators. It removes the lower-through-strand
summands and makes simple projections explicit.

\item[Lagrangian algebra.]A connected commutative separable algebra
whose squared dimension equals the global dimension of the
ambient modular category; its local module category is trivial.

\item[Marginal by dimension.]A scalar of dimension two. Exact
marginality additionally requires vanishing beta functions and
compatibility among deformations.

\item[Modular category.]A braided ribbon fusion category with
nondegenerate braiding; equivalently in this setting, its
normalized $S$ matrix is invertible. Positivity is a further
property when a unitary structure is specified.

\item[Module category.]A category on which a fusion category acts,
with coherent associativity maps. It encodes boundary and gapped
vacuum patterns unavailable from the fusion ring alone.

\item[Morita equivalence.]Equivalence through an invertible
bimodule-category relation. Morita-equivalent fusion categories
have equivalent centers but need not have the same simple objects
or the same chain Hamiltonian.

\item[Near-group rules.]Fusion rules with a group of invertible
simples and only one noninvertible simple. Generalized Haagerup
rules instead have one noninvertible simple for each group
element.

\item[Normal convergence.]Uniform absolute convergence on
compact subsets of the relevant complex parameter domain.
It permits the analytic operations used with the trace
functions in the stated domains.

\item[Ordinary module.]A module with diagonalizable energy,
finite-dimensional eigenspaces, and lower truncation along
each integer-shift weight string. A simple ordinary module
here has one lowest-weight grading; a direct sum can have several.

\item[Perron eigenvector.]A positive eigenvector of a nonnegative
matrix at its spectral radius. For a matrix with strongly connected
directed graph it is unique up to scale; primitivity additionally
excludes other eigenvalues of the same modulus. Such vectors supply
positive weights in fusion graphs and classical transfer matrices.

\item[Planar algebra.]Vector spaces of labeled boxes equipped
with multilinear operations specified by planar string diagrams.
Substituting diagrams agrees with composition of operations.
A positive subfactor planar algebra reconstructs tensor categories
by taking sums and splitting projections.

\item[Primary versus quasiprimary.]A Virasoro primary is killed by
all positive Virasoro modes; a quasiprimary need only be killed by
$L_1$ and have definite weight.

\item[Primitive transfer operator.]A finite-dimensional completely
positive transfer map for which, after normalization and restriction
to its supported algebra, the eigenvalue one is simple, all other
eigenvalues have modulus less than one, and the left and right
fixed metrics are strictly positive. Equivalently in this setting,
a suitable positive power sends every nonzero positive operator
to a strictly positive operator. A symmetry cat can fail this
property even when its pure components are well behaved.

\item[Projective object.]An object $P$ for which any morphism
$P\to Y$ lifts through every surjection $X\to Y$.
Projectivity of the vacuum is a central step in the
regularity argument; right exactness of a functor is a
different property.

\item[Pseudounitary category.]A fusion category whose
categorical and Frobenius--Perron global dimensions agree,
with the corresponding canonical positive spherical dimensions.
The definition does not supply a positive adjoint on morphisms.

\item[Q-system.]A positive separable Frobenius algebra object in a
unitary tensor category, in the appropriate normalization. It is
the algebraic datum of a finite-index extension or a module
construction.

\item[Ramified prime.]For a number field, a rational prime
dividing its field discriminant. Ramification persists in a larger
number field, which lets it obstruct a proposed field inclusion.
For $\Q(\zeta_N)$, only primes dividing $N$ can ramify.

\item[Regularity.]Complete decomposition of weak modules into
ordinary simple modules, in the VOA sense used in the construction.

\item[Right exact functor.]A functor preserving cokernels,
and therefore sequences $X\to Y\to Z\to0$ in the relevant
abelian setting. It need not preserve kernels or split extensions.

\item[Schmidt decomposition.]The singular-value decomposition
of a pure-state coefficient matrix across a bipartition.
Its squared singular values are the reduced density-matrix
eigenvalues and determine the bipartite entropy.

\item[Sewing.]Compatibility of correlation functions under cutting,
inserting a complete set of states, and gluing surfaces. Sphere
crossing and torus modular invariance are related parts of this
larger requirement.

\item[Smith normal form.]An integer matrix diagonalized using
invertible integer row and column operations, with successive
nonzero diagonal entries dividing the next. Those entries
describe the cyclic factors of its presented abelian group.

\item[Spherical structure.]A coherent choice of duality traces
for which left and right categorical traces agree. It permits
isotopy-invariant loop evaluations and categorical dimensions.

\item[Split rational form.]A tensor category defined over a
smaller field with its simple objects and scalar endomorphism
algebras already split there. Extending scalars gives the complex
category. A field of dimensions alone need not be such a field
of definition.

\item[Strong generation.]Spanning by derivatives and normal
products of a set of fields. It is stronger than arbitrary-mode
vertex generation.

\item[Strongly rational VOA.]A simple, self-contragredient,
CFT-type VOA that is rational and $C_2$-cofinite, with the
module-theoretic conventions stated in the regularity chapter.

\item[Tensorator.]The coherent isomorphisms comparing the image
of a tensor product with the tensor product of images under
a monoidal functor. A linear equivalence does not supply them.

\item[Tube algebra.]A finite algebra of annular line networks whose
simple representations encode center charge sectors of a fusion
category. Its finite representation dimension is not a quantum
dimension or the dimension of the dynamical CFT sector. It counts
copies of that dynamical sector and therefore multiplies its
state degeneracies.

\item[Twined versus twisted.]A twined trace inserts a line operator
in an ordinary trace; a twisted trace counts states in a Hilbert
space with a defect crossing the spatial circle.

\item[Unitary fusion category.]A fusion category with compatible
positive adjoints on morphisms and unitary structural maps.
Positive Frobenius--Perron dimensions alone do not imply
this structure.

\item[Weakly integral.]A fusion category whose global
Frobenius--Perron dimension is an integer. Its simple dimensions
have integer squares but need not be integers: the Ising
dimensions $1,1,\sqrt2$ give global dimension four.
The square condition concerns simple objects, not their arbitrary sums.

\item[Weak module.]A VOA module satisfying the field and
Jacobi identities without a preassigned ordinary energy grading.
Regularity supplies a decomposition into ordinary simple modules;
it is not part of the weak-module definition.

\item[Zhu algebra.]An associative algebra built from a VOA
whose representations describe the degree-zero action on
module top spaces. Its use here relates complete component
tops to the spectrum of partial conformal energies.

\end{description}
